\documentclass{article}
\usepackage[margin=1in]{geometry}
\usepackage{amsmath,amssymb,amsthm}

\newcommand{\Actor}{\mathsf{Actor}}
\newcommand{\Checkpoint}{\mathsf{Checkpoint}}
\newcommand{\Action}{\mathsf{Action}}
\newcommand{\State}{\Sigma}
\newcommand{\Obs}{\mathsf{Obs}}
\newcommand{\Model}{\mathsf{ModelState}}
\newcommand{\Input}{\mathsf{Input}}
\newcommand{\Mapping}{\mathsf{Mapping}}
\newcommand{\Mode}{\mathsf{Mode}}
\newcommand{\Phase}{\mathsf{Phase}}
\newcommand{\Fault}{\mathsf{Fault}}
\newcommand{\Event}{\mathsf{Event}}
\newcommand{\Receipt}{\mathsf{Receipt}}
\newcommand{\Config}{\mathsf{Config}}
\newcommand{\Parked}{\mathsf{parked}}
\newcommand{\Running}{\mathsf{running}}
\newcommand{\Done}{\mathsf{done}}
\newcommand{\Quiet}{\mathsf{quiet}}
\newcommand{\Run}{\mathsf{run}}
\newcommand{\CheckPhase}{\mathsf{check}}
\newcommand{\Failed}{\mathsf{failed}}
\newcommand{\Succeeded}{\mathsf{succeeded}}
\newcommand{\CPDone}{\mathsf{\$done}}
\newcommand{\Internal}{\mathsf{internal}}
\newcommand{\Stop}{\mathsf{stop}}
\newcommand{\project}{\mathsf{project}}
\newcommand{\ScheduleOf}{\mathsf{schedule}}
\newcommand{\observe}{\mathsf{observe}}
\newcommand{\matches}{\approx}
\newcommand{\ReplayValid}{\mathsf{ReplayValid}}
\newcommand{\Merge}{\mathsf{Merge}}

\title{A Type-Theoretic Model of DPM Replay}
\author{}
\date{}

\begin{document}
\maketitle

\section{Scope}

This paper gives a proposed formal model of the core deterministic scheduling
(DPM) replay mechanism in Mirrors. It covers admitted actor/checkpoint
schedules, cooperative execution of actual application workers, observation of
actual system-under-test (SUT) states, and comparison with a model trace. It
uses the restricted profile in which every scheduled checkpoint interval,
including $\CPDone$, is the image of one fully instantiated model action and
has one subsequent observation. Spectator initialization and richer native or
multi-phase mappings are outside this profile. The paper does not formalize
artifact identities, the command-line interface, reduction, native process
isolation, installation, timeouts, cancellation, actor or observer exceptions,
undeclared workers, cleanup failures, or failure recovery.

The central separation is:
\[
  \begin{gathered}
    \text{DPM requires and verifies schedule fidelity on successful replay,}\\
    \text{while Mirrors checks whether actual observations correspond to
      expected model states.}
  \end{gathered}
\]
Expected model states are oracles, not inputs used to initialize or steer the
SUT.

\section{PFPL basis and notation discipline}

We assume as already defined the four transition-system judgment forms from
PFPL, second edition, Section~5.1 \cite{pfpl}:
\[
  s\ \mathsf{state},\qquad
  s\ \mathsf{initial},\qquad
  s\ \mathsf{final},\qquad
  s\longmapsto s'.
\]
We also assume PFPL's iterated judgment $s\longmapsto^{*}s'$. We use its
labeled process-dynamics judgment (Section~39.2) as well. The DPM
configuration dynamics below uses the unlabeled judgment $C\longmapsto C'$.
For SUT workers it uses the instance
\[
  \sigma\xrightarrow{\lambda}_{a}\sigma'.
\]
The subscript $a$ selects the transition system of actor $a$, and $\lambda$ is
the paper-specific worker label defined in Section~\ref{sec:app-parameters}.
No general PFPL judgment is redefined here. In our instantiation, finality
alone does not assert replay success because both successful and failed
configurations are final.

PFPL notation and ordinary mathematical logic are not redefined. Every symbol
introduced specifically for this paper is defined below. Application-specific
objects are defined as explicit parameters with stated types and meanings;
none is an undocumented global assumption.

\section{Basic types and operators}

\subsection{Paper-specific types}

\begin{center}
\begin{tabular}{ll}
  Symbol & Definition\\ \hline
  $\Actor$ & finite type of declared actor identifiers\\
  $\Checkpoint$ & finite type of checkpoint identifiers\\
  $\Action$ & type of fully instantiated model actions, including parameters\\
  $\State$ & actual SUT states, including shared state and worker continuations\\
  $\Obs$ & type of observations exposed to Mirrors\\
  $\Model$ & type of model states\\
  $\Input$ & type of recorded actual inputs\\
  $\Mapping$ & type of model-action to actor/checkpoint mappings\\
  $\Mode$ & actor-mode type, defined in Section~\ref{sec:configuration-statics}\\
  $\Phase$ & coordinator-phase type, defined in Section~\ref{sec:configuration-statics}\\
  $\Fault$ & terminal-fault type, defined in Section~\ref{sec:configuration-statics}\\
  $\Event$ & receipt-event type, defined in Section~\ref{sec:configuration-statics}\\
  $\Receipt$ & $\mathsf{List}(\Event)$, the type of retained receipts\\
  $\Config$ & replay-configuration type, defined in Section~\ref{sec:configuration-statics}
\end{tabular}
\end{center}

The special checkpoint $\CPDone\in\Checkpoint$ denotes actual worker
completion, including its relevant effects. The letters $a,b$ range over
$\Actor$, $\kappa$ over $\Checkpoint$, $\alpha$ over $\Action$, $\sigma$ over
$\State$, $o$ over $\Obs$, $s$ over $\Model$, and $I$ over $\Input$.

\subsection{Lists and finite maps}

For a type $X$, $\mathsf{List}(X)$ is the type of finite lists over $X$. We
write $\epsilon$ for the empty list, $x::xs$ for cons, and $|xs|$ for length.
When $r$ is a receipt and $\eta$ is one event, $r\cdot\eta$ appends $\eta$ at
the end of $r$. These conventions fix the use of list notation in this paper;
they are not additional transition judgments.

The symbol $\rightharpoonup$ denotes a finite partial map, and
$\operatorname{dom}(D)$ denotes its domain. If $\rho$ is a finite map,
$\rho[a\mapsto m]$ is the map identical to $\rho$ except that its value at $a$
is $m$. The declared actor set is
\[
  \mathcal A=\operatorname{dom}(D).
\]
The notation $\rho:\mathcal A\to\Mode$ means that $\rho$ is total on the finite
set $\mathcal A$.

\subsection{Schedules, traces, and configurations}

A schedule step is $(a,\kappa)\in\Actor\times\Checkpoint$. A schedule $Q$ is a
list of such steps. The declaration
\[
  D:\Actor\rightharpoonup\mathsf{List}(\Checkpoint)
\]
assigns a local checkpoint chain to each declared actor. The full admitted
schedule is written $Q_0$; $Q$ denotes its unexecuted suffix in a configuration.

A model trace with $n$ transitions is
\[
  \tau=s_0,\alpha_0,s_1,\ldots,\alpha_{n-1},s_n,
\]
where each $s_j$ is a model state and each $\alpha_j$ is a fully instantiated
model action. Thus two actions with different actor or input parameters are
different elements of $\Action$. A mapping is a function
\[
  \mu:\Action\to\Actor\times\Checkpoint
\]
defined at least on the actions occurring in $\tau$. Its schedule interpretation
is defined by
\[
  \ScheduleOf_\mu(\tau)
  =(\mu(\alpha_0),\ldots,\mu(\alpha_{n-1})).
\]
Consequently $|Q_0|=n$. The restricted profile requires the terminal
$\CPDone$ steps needed by the actor chains to occur among these $n$ mapped
actions; it does not append unmodeled teardown steps.

The configuration variables have the following meanings:
\begin{center}
\begin{tabular}{ll}
  Symbol & Definition\\ \hline
  $\rho$ & finite map from declared actors to their modes\\
  $\sigma$ & current actual SUT state\\
  $\pi$ & current coordinator phase\\
  $i$ & index of the next model state to compare, $0\le i\le n+1$\\
  $r$ & receipt accumulated so far\\
  $\eta$ & one receipt event\\
  $e$ & one terminal fault\\
  $\lambda$ & one labeled SUT-worker action
\end{tabular}
\end{center}

\section{Application-supplied parameters}
\label{sec:app-parameters}

An instantiation of the model must supply the following parameters:
\[
\begin{array}{ll}
  \operatorname{Init}\subseteq\Input\times\State
    & \text{actual SUT initialization from recorded input},\\
  \observe:\State\to\Obs
    & \text{actual-state observation function},\\
  {\matches}\subseteq\Obs\times\Model
    & \text{Mirrors observation/model correspondence},\\
  \operatorname{CleanupOk}\subseteq\State
    & \text{states with required joins and teardown confirmed}.
\end{array}
\]
Thus $\operatorname{Init}(I,\sigma)$ means that actual input $I$ can construct
actual SUT state $\sigma$; it has no model-state argument. The judgment
$o\matches s$ means that observation $o$ corresponds to model state $s$ under
the supplied Mirrors oracle. The predicate $\operatorname{CleanupOk}(\sigma)$
does not assert model correspondence.

For each $a\in\mathcal A$, the application also supplies a PFPL-style labeled
transition relation $\sigma\xrightarrow{\lambda}_{a}\sigma'$. Its labels are
defined inductively by
\[
  \lambda::=\Internal\mid\mathsf{arrive}(\kappa)\mid\Stop .
\]
The label $\Internal$ is an actual application step that does not report a
checkpoint; $\mathsf{arrive}(\kappa)$ reports cooperative arrival at $\kappa$
and parks the worker while retaining its continuation unless
$\kappa=\CPDone$; and $\Stop$ reports worker return without such an arrival.
The expected model trace is not a parameter of this SUT transition relation.

\section{Schedule statics}

The projection $\project_a(Q)$ retains exactly the steps of actor $a$:
\[
\begin{aligned}
  \project_a(\epsilon)&=\epsilon,\\
  \project_a((b,\kappa)::Q)&=
  \begin{cases}
    (a,\kappa)::\project_a(Q)&a=b,\\
    \project_a(Q)&a\ne b.
  \end{cases}
\end{aligned}
\]
The set $\operatorname{actors}(Q)$ contains every actor occurring in $Q$. A
chain $c$ satisfies $\operatorname{wellChain}(c)$ iff it is nonempty, its last
checkpoint is $\CPDone$, and no earlier checkpoint is $\CPDone$.

The paper-specific admission judgment
$D;\tau;\mu\vdash Q_0\ \mathsf{admissible}$ is defined inductively by
\[
\frac{
  \begin{array}{c}
    Q_0=\ScheduleOf_\mu(\tau)\qquad
    \operatorname{actors}(Q_0)=\mathcal A\\
    \forall a\in\mathcal A.\
      \operatorname{wellChain}(D(a))\ \land\
      \project_a(Q_0)=D(a)
  \end{array}
}{
  D;\tau;\mu\vdash Q_0\ \mathsf{admissible}
}
\]
Every actor therefore preserves its local chain, while $Q_0$ selects one
interleaving derived from the model trace. Admission does not quantify over
alternative interleavings.

\section{Configuration statics}
\label{sec:configuration-statics}

The metavariables $m$, $\pi$, $e$, and $\eta$ range respectively over
$\Mode$, $\Phase$, $\Fault$, and $\Event$. These four inductive types have the
following constructors:
\[
\begin{array}{rcl}
m&::=&\Parked\mid\Running\mid\Done,\\[2pt]
\pi&::=&\Quiet\mid\Run(a,\kappa)\mid\CheckPhase(j)
       \mid\Failed(e)\mid\Succeeded,\\[2pt]
e&::=&\mathsf{unexpected}(a,\kappa,b,\kappa')
       \mid\mathsf{premature}(a)
       \mid\mathsf{modelMismatch}(j,o,s),\\[2pt]
\eta&::=&\mathsf{permit}(a,\kappa)
       \mid\mathsf{arrival}(a,\kappa)
       \mid\mathsf{observation}(j,o)
       \mid\mathsf{fault}(e)
       \mid\mathsf{cleanup}.
\end{array}
\]
In $\CheckPhase(j)$ and $\mathsf{modelMismatch}(j,o,s)$, $j$ is a model-state
index satisfying $0\le j\le n$. The constructors $\Run$ and $\CheckPhase$ are
coordinator phases, not SUT-program values. In
$\mathsf{unexpected}(a,\kappa,b,\kappa')$, the pair $(a,\kappa)$ is the
expected actor/checkpoint step and $(b,\kappa')$ is the actual arrival.
The receipt grammar is semantic rather than a raw SDK schema: in particular,
$\mathsf{permit}(a,\kappa)$ records a grant whose expected destination is
$\kappa$. Concrete SDK receipts may instead encode actor-local permit origins
and schedule-step indices.

A configuration is
\[
  C=(Q,\rho,\sigma,\pi,i,r)\in\Config.
\]
Following PFPL's $\mathsf{state}$ judgment, configuration formation is defined
by
\[
\frac{
  \begin{array}{c}
    Q\in\mathsf{List}(\Actor\times\Checkpoint)\qquad
    \rho:\mathcal A\to\Mode\qquad
    \sigma:\State\\
    \pi:\Phase\qquad 0\le i\le n+1\qquad r:\Receipt
  \end{array}
}{
  (Q,\rho,\sigma,\pi,i,r)\ \mathsf{state}
}
\]
This is syntactic state formation. Execution invariants such as mutual
exclusion of running actors are proved separately rather than assumed here.

The initial-state judgment is defined by
\[
\frac{
  \begin{array}{c}
    D;\tau;\mu\vdash Q_0\ \mathsf{admissible}\qquad
    \operatorname{Init}(I,\sigma_0)\\
    \rho_0=\{(a,\Parked)\mid a\in\mathcal A\}
  \end{array}
}{
  (Q_0,\rho_0,\sigma_0,\CheckPhase(0),0,\epsilon)\ \mathsf{initial}
}
\]
The notation $\{(a,\Parked)\mid a\in\mathcal A\}$ denotes the finite map that
assigns $\Parked$ to every declared actor. The initial observation is checked
by the ordinary comparison dynamics below; it is not assumed to match.

\section{Structural dynamics of replay}

All rules below define the single PFPL transition judgment $C\longmapsto C'$
for the fixed admitted replay instance $(D,Q_0,\tau,I,\mu)$.

\subsection*{Initial or subsequent observation matches}
\[
\frac{
  j=i\qquad o=\observe(\sigma)\qquad o\matches s_j
}{
  (Q,\rho,\sigma,\CheckPhase(j),i,r)
  \longmapsto
  (Q,\rho,\sigma,\Quiet,i+1,r\cdot\mathsf{observation}(j,o))
}
\]

\subsection*{Observation mismatch}
\[
\frac{
  j=i\qquad o=\observe(\sigma)\qquad \neg(o\matches s_j)
}{
  (Q,\rho,\sigma,\CheckPhase(j),i,r)
  \longmapsto
  (Q,\rho,\sigma,\Failed(\mathsf{modelMismatch}(j,o,s_j)),i,
    r\cdot\mathsf{fault}(\mathsf{modelMismatch}(j,o,s_j)))
}
\]

\subsection*{Permit one actor}
\[
\frac{
  \rho(a)=\Parked
}{
  ((a,\kappa)::Q,\rho,\sigma,\Quiet,i,r)
  \longmapsto
  ((a,\kappa)::Q,\rho[a\mapsto\Running],\sigma,\Run(a,\kappa),i,
    r\cdot\mathsf{permit}(a,\kappa))
}
\]

\subsection*{Actual internal SUT step}
\[
\frac{
  a\in\mathcal A\qquad \rho(a)=\Running\qquad
  \sigma\xrightarrow{\Internal}_{a}\sigma'
}{
  (Q,\rho,\sigma,\Run(a,\kappa),i,r)
  \longmapsto
  (Q,\rho,\sigma',\Run(a,\kappa),i,r)
}
\]
This rule exposes arbitrary real execution inside a checkpoint interval; it
does not treat that interval as one instruction-level atomic operation.

\subsection*{Expected actual arrival}
\[
\frac{
  \begin{array}{c}
    a\in\mathcal A\qquad \rho(a)=\Running\qquad 1\le i\le n\\
    \sigma\xrightarrow{\mathsf{arrive}(\kappa)}_{a}\sigma'
  \end{array}\qquad
  m=\begin{cases}
    \Done&\kappa=\CPDone,\\
    \Parked&\text{otherwise}
  \end{cases}
}{
  ((a,\kappa)::Q,\rho,\sigma,\Run(a,\kappa),i,r)
  \longmapsto
  (Q,\rho[a\mapsto m],\sigma',\CheckPhase(i),i,
    r\cdot\mathsf{arrival}(a,\kappa))
}
\]
The index guard ensures that the resulting $\CheckPhase(i)$ is well formed.
Reachable running states satisfy it because the initial observation precedes
the first permit and each arrival has one corresponding model transition.

\subsection*{Unexpected actual arrival}
\[
\frac{
  \begin{array}{c}
    a\in\mathcal A\qquad \rho(a)=\Running\qquad b\in\mathcal A\\
    \sigma\xrightarrow{\mathsf{arrive}(\kappa')}_{b}\sigma'\qquad
    (b,\kappa')\ne(a,\kappa)\\
    e=\mathsf{unexpected}(a,\kappa,b,\kappa')
  \end{array}
}{
  ((a,\kappa)::Q,\rho,\sigma,\Run(a,\kappa),i,r)
  \longmapsto
  ((a,\kappa)::Q,\rho,\sigma',\Failed(e),i,r\cdot\mathsf{fault}(e))
}
\]
There is no rule that reorders the schedule or repairs an unexpected arrival.

\subsection*{Premature worker return}
\[
\frac{
  a\in\mathcal A\qquad \rho(a)=\Running\qquad
  \sigma\xrightarrow{\Stop}_{a}\sigma'\qquad
  e=\mathsf{premature}(a)
}{
  ((a,\kappa)::Q,\rho,\sigma,\Run(a,\kappa),i,r)
  \longmapsto
  ((a,\kappa)::Q,\rho,\sigma',\Failed(e),i,r\cdot\mathsf{fault}(e))
}
\]

\subsection*{Successful cleanup and completion}
\[
\frac{
  \forall a\in\mathcal A.\,\rho(a)=\Done\qquad
  \operatorname{CleanupOk}(\sigma)
}{
  (\epsilon,\rho,\sigma,\Quiet,n+1,r)
  \longmapsto
  (\epsilon,\rho,\sigma,\Succeeded,n+1,r\cdot\mathsf{cleanup})
}
\]

\subsection*{Final states}
\[
\frac{
  C=(Q,\rho,\sigma,\Succeeded,i,r)\qquad C\ \mathsf{state}
}{
  C\ \mathsf{final}
}
\qquad
\frac{
  C=(Q,\rho,\sigma,\Failed(e),i,r)\qquad C\ \mathsf{state}
}{
  C\ \mathsf{final}
}
\]
No transition rule has $\Succeeded$ or $\Failed(e)$ as its source phase, so
these states satisfy PFPL's requirement that final states have no outgoing
transition. A failed final state is not a successful replay. Moreover, a
failed final state is terminal only for this replay machine: it does not
assert that an actor whose recorded mode is still $\Running$ has stopped or
been reclaimed. Cancellation, cleanup retries, and failure recovery are
outside the stated scope.

\section{Replay correctness}

Successful replay is the paper-specific predicate
\[
\begin{aligned}
  \ReplayValid(D,Q_0,\tau,I,\mu)
  \iff{}&\exists \sigma_0,\rho_s,\sigma_s,r_s.\
    (Q_0,\rho_0,\sigma_0,\CheckPhase(0),0,\epsilon)\ \mathsf{initial}\\
    &\land (Q_0,\rho_0,\sigma_0,\CheckPhase(0),0,\epsilon)
      \longmapsto^{*}
      (\epsilon,\rho_s,\sigma_s,\Succeeded,n+1,r_s)\\
    &\land (\epsilon,\rho_s,\sigma_s,\Succeeded,n+1,r_s)\ \mathsf{final}
\end{aligned}
\]

A configuration $C$ is \emph{reachable} iff there exists an initial
configuration $C_0$ such that $C_0\longmapsto^{*}C$. Every initial
configuration satisfies $C_0\ \mathsf{state}$ by the formation and admission
rules.

The following properties follow by rule induction on $C\longmapsto C'$ or on
its PFPL iteration. The induction simultaneously carries the consumed schedule
prefix, the phase/mode invariant, and the next-model-index invariant; the final
state alone is not a sufficient induction hypothesis:
\begin{enumerate}
\item \textbf{Preservation of state formation.}
  If $C\ \mathsf{state}$ and $C\longmapsto C'$, then $C'\ \mathsf{state}$.

\item \textbf{Schedule fidelity.}
  In a successful replay receipt, the sequence of steps carried by permit
  events is exactly $Q_0$, and the sequence carried by actual-arrival events is
  also exactly $Q_0$. Each admitted step therefore has both its scheduled
  permit and its expected actual arrival; observation and cleanup events are
  not schedule steps.

\item \textbf{Coordinator mode invariant.}
  In every reachable configuration, at most one declared actor has coordinator
  mode $\Running$. In every nonfinal reachable configuration, that mode occurs
  only in a $\Run(a,\kappa)$ phase. A failed final configuration may retain
  the $\Running$ mode to record an unreclaimed actor. An actual arrival by a
  different declared actor is rejected by the unexpected-arrival rule rather
  than being hidden by this coordinator invariant.

\item \textbf{Quiescent observation.}
  No observation occurs while a declared actor is $\Running$. Observations are
  made from actual SUT states initially or after the selected actor has parked
  or finished.

\item \textbf{Observation correspondence.}
  Successful replay produces actual observations $o_0,\ldots,o_n$ such that
  \[
    \forall j\in\{0,\ldots,n\}.\quad o_j\matches s_j.
  \]

\item \textbf{Oracle separation.}
  Within the dynamics, model states are used only in comparison rules and
  resulting fault evidence. They are not arguments of $\operatorname{Init}$ or
  of any labeled SUT transition. Model actions still influence the schedule
  through $\mu$. This is a syntactic interface separation, not a proof that an
  arbitrary application adapter is faithful; an instantiation must supply
  initialization and SUT relations that do not secretly consult model states.

\item \textbf{Single-path scope.}
  $\ReplayValid$ establishes one witnessed conforming execution along $Q_0$.
  It does not assert that every resolution of SUT nondeterminism conforms, or
  quantify over other schedules, unobserved SUT states, undeclared actors,
  weak-memory executions, or preemption inside checkpoint intervals.
\end{enumerate}
No progress or termination theorem is claimed. Initialization may have no
result, a worker may diverge through $\Internal$ steps, required cleanup may
be unconfirmed, and a nonfinal configuration may be stuck. $\ReplayValid$ is
an existential statement about a complete successful replay, not a guarantee
that every admitted replay terminates.

\section{Optional finite-interleaving relation}

For lists $x,y,z$ of schedule steps, the order-preserving merges used by finite
exploration are defined by the paper-specific inductive judgment
$\Merge(x,y,z)$:
\[
\frac{}{\Merge(\epsilon,\epsilon,\epsilon)}
\]
\[
\frac{\Merge(xs,ys,zs)}
     {\Merge(x::xs,ys,x::zs)}
\qquad
\frac{\Merge(xs,ys,zs)}
     {\Merge(xs,y::ys,y::zs)}
\]
The variables $xs$, $ys$, and $zs$ denote the respective list tails. This
relation describes candidate schedules only. Complete coverage of this finite
merge space additionally requires a known, untruncated denominator, completed
execution of every required schedule, and confirmed cleanup. An early mismatch
does not complete the unexecuted suffix. Even complete finite merge coverage
does not by itself establish model comparison or universal concurrent-state
coverage.

\begin{thebibliography}{1}
\bibitem{pfpl}
Robert Harper.
\emph{Practical Foundations for Programming Languages}.
Second edition, Cambridge University Press, 2016.
Sections~5.1 and~39.2.
\end{thebibliography}

\end{document}
